%%%%%%%%%%%%%%%%%%%%%
% Documento maestro %
%%%%%%%%%%%%%%%%%%%%%
\documentclass{fime}

%%%%%%%%%%%%%%%%%%%%%%%%%%%%%%%%%%%%%%%%%%%
% Cargando paquetes y definiendo opciones %
%%%%%%%%%%%%%%%%%%%%%%%%%%%%%%%%%%%%%%%%%%%
% Aquí puedes cargar los paquetes que vas a usar. La clase
% fime ya incluye babel, inputenc, graphicx y los de la AMS.
% Cargar un paquete está a tu libertad (y responsabilidad).
\usepackage{hyperref}
    \hypersetup{breaklinks=true,colorlinks=true,
        linkcolor=black,citecolor=black,urlcolor=black}
\usepackage{datetime} % monthname

%%%%%%%%%%%%%%%%%%%%%
% Definiendo campos %
%%%%%%%%%%%%%%%%%%%%%
\def\titulo{Optimal Location of Car Wreck Adjusters}
\def\autor{Luis Alberto Maltos Ortega}
\def\matricula{1390200}
\def\grado{Maestría en Ciencias en Ingenie\'ia de Sistemas}
% En caso de que el grado tenga orientación o especialidad llenar el siguiente
% campo dejando un ESPACIO INICIAL, en caso contrario, dejar vacío
\def\orientacion{ }
\def\fecha{\monthname,  \the\year} % Coloca el mes con mayúscula inicial

\def\asesor{Dr. Roger Z. R\'ios}
\def\revisorA{Mar\'ia Ang\'elica Salazar Aguilar}
\def\revisorB{Mar\'ia Guadalupe Villareal Marroqu\'in}
% En el caso de que tu tesis sea de doctorado activa la variable cambiándola a \doctoradotrue
% y define tus otros dos revisores
\newif\ifdoctorado\doctoradofalse
\def\revisorC{Nombre del revisor C}
\def\revisorD{Nombre del revisor D}
% El visto bueno siempre va
\def\vobo{Dr. Simón Martínez Martínez}

%%%%%%%%%%%%%%%%%%%%%%%
% Inicia el documento %
%%%%%%%%%%%%%%%%%%%%%%%
\begin{document}

\frontmatter
\pagestyle{main}

%%% Incluye PortillasM si tu tesis es de Maestría
%%% y PortillasD si es de doctorado.
% Portadas (Maestría)

\def\uanl{Universidad Autónoma de Nuevo León}
\def\fime{Facultad de Ingeniería Mecánica y Eléctrica}
\def\depg{Subdirección de Estudios de Posgrado}
\def\snnl{San Nicolás de los Garza, Nuevo León}

%%%%%%%%%%%%%%%%%%%%%%%%
% Primer portada: UANL %
%%%%%%%%%%%%%%%%%%%%%%%%
\thispagestyle{empty}

\begin{scshape}
\begin{center}
	{\Large \uanl} \\[5mm]
	{\large \fime} \\[5mm]
	{\large \depg}
	\vskip 15mm
	\includegraphics[height=55mm]{uanl}
	\vskip 12mm
	\begin{tabular}{p{11cm}}
		\centering
		{\large \titulo}
	\end{tabular}
	\vskip 7mm
	{por}\\[7mm]
	{\large \autor}\\[7mm]
	{como requisito parcial para obtener el grado de}\\[3mm]
	\MakeUppercase{\grado}\\
	\orientacion
	\vfill
	\fecha
\end{center}
\end{scshape}

%%%%%%%%%%%%%%%%%%%%%%%%%
% Segunda portada: FIME %
%%%%%%%%%%%%%%%%%%%%%%%%%
\newpage
\thispagestyle{empty}

\begin{scshape}
\begin{center}
	{\Large \uanl} \\[5mm]
	{\large \fime} \\[5mm]
	{\large \depg}
	\vskip 16mm
	\includegraphics[height=50mm]{fime}
	\vskip 16mm
	\begin{tabular}{p{11cm}}
		\centering
		{\large \titulo}
	\end{tabular}
	\vskip 7mm
	{por}\\[7mm]
	{\large \autor}\\[7mm]
	{como requisito parcial para obtener el grado de}\\[3mm]
	\MakeUppercase{\grado}\\
	\orientacion
	\vfill
	\fecha
\end{center}
\end{scshape}

%%%%%%%%%%%%%%%%%%%%%%%%%%%%%
% Carta del comité de tesis %
%%%%%%%%%%%%%%%%%%%%%%%%%%%%%
\newpage
\thispagestyle{empty}
\enlargethispage{5mm}

\begin{center}
{\bf \large \uanl} \\
{\bf \fime} \\
{\bf \depg}
\end{center}
\vskip 4mm

Los miembros del Comité de Tesis recomendamos que la Tesis <<\titulo>>, realizada por el alumno \autor, con número de matrícula \matricula, sea aceptada para su defensa como requisito parcial para obtener el grado de \grado\orientacion.
\ifdoctorado{\vskip 10mm}\else{\vskip 8mm}\fi

\begin{center}
El Comité de Tesis\\
\ifdoctorado{\vskip 15mm}\else{\vskip 25mm}\fi

\ifdoctorado{%%%
\begin{tabular}{p{37mm}p{21mm}p{12mm}p{21mm}p{37mm}}
	\cline{2-4}
	& \multicolumn{3}{c}{\asesor} & \\
	& \multicolumn{3}{c}{Asesor}  & \\[15mm]
	\cline{1-2} \cline{4-5}
	\multicolumn{2}{c}{\revisorA} & & \multicolumn{2}{c}{\revisorB} \\
	\multicolumn{2}{c}{Revisor}   & & \multicolumn{2}{c}{Revisor}   \\[17mm]
	\cline{1-2} \cline{4-5}
	\multicolumn{2}{c}{\revisorC} & & \multicolumn{2}{c}{\revisorD} \\
	\multicolumn{2}{c}{Revisor}   & & \multicolumn{2}{c}{Revisor}   \\[2mm]
	& \multicolumn{3}{c}{Vo. Bo.} & \\[14mm]
	\cline{2-4}
	& \multicolumn{3}{c}{\vobo}   & \\
	& \multicolumn{3}{c}{\depg}   & \\ &&&&
\end{tabular}
}\else{%%%
\begin{tabular}{p{37mm}p{21mm}p{12mm}p{21mm}p{37mm}}
	\cline{2-4}
	& \multicolumn{3}{c}{\asesor} & \\
	& \multicolumn{3}{c}{Asesor}  & \\[19mm]
	\cline{1-2} \cline{4-5}
	\multicolumn{2}{c}{\revisorA} & & \multicolumn{2}{c}{\revisorB} \\
	\multicolumn{2}{c}{Revisor}   & & \multicolumn{2}{c}{Revisor}   \\[2mm]
	& \multicolumn{3}{c}{Vo. Bo.} & \\[17mm]
	\cline{2-4}
	& \multicolumn{3}{c}{\vobo}   & \\
	& \multicolumn{3}{c}{\depg}   & \\ &&&&
\end{tabular}
}\fi%%%

\vfill

\snnl, \MakeLowercase{\fecha}

\end{center}

% Dedicatoria

\thispagestyle{empty}
\vspace*{17mm}

\begin{flushright}
\begin{itshape}

Aquí puedes poner tu dedicatoria\\
si es que tienes una.\bigskip\bigskip

Si no tienes una, puedes borrar\\
la línea \verb+% Dedicatoria

\thispagestyle{empty}
\vspace*{17mm}

\begin{flushright}
\begin{itshape}

Aquí puedes poner tu dedicatoria\\
si es que tienes una.\bigskip\bigskip

Si no tienes una, puedes borrar\\
la línea \verb+% Dedicatoria

\thispagestyle{empty}
\vspace*{17mm}

\begin{flushright}
\begin{itshape}

Aquí puedes poner tu dedicatoria\\
si es que tienes una.\bigskip\bigskip

Si no tienes una, puedes borrar\\
la línea \verb+\include{Dedicatoria}+ en el\\
archivo \texttt{MiTesis.tex} pues no es obligatoria.

\end{itshape}
\end{flushright}

+ en el\\
archivo \texttt{MiTesis.tex} pues no es obligatoria.

\end{itshape}
\end{flushright}

+ en el\\
archivo \texttt{MiTesis.tex} pues no es obligatoria.

\end{itshape}
\end{flushright}



\tableofcontents
\listoffigures
\listoftables

%Agradecimientos

\chapter{Agradecimientos}
\markboth{Agradecimientos}{}

Aquí puedes poner tus agradecimientos. (No olvides agradecer a tu comité de tesis, a tus profesores, a la facultad y a CONACyT en caso de que hallas sido beneficiado con una beca).

%Resumen

\chapter{Resumen}
\markboth{Resumen}{}

{\setlength{\leftskip}{10mm}
\setlength{\parindent}{-10mm}

\autor.

Candidato para obtener el grado de \grado\orientacion.

\uanl.

\fime.

Título del estudio: \textsc{\titulo}.

\noindent Número de páginas: \pageref*{lastpage}.}

%%% Comienza a llenar aquí
\paragraph{Objetivos y método de estudio:}
Aquí debes poner tus objetivos y métodos de estudio. (Este es el formato).

\paragraph{Contribuciones y conlusiones:}
Y aquí tus contribuciones y conclusiones. (También es parte del formato).

\bigskip\noindent\begin{tabular}{lc}
\vspace*{-2mm}\hspace*{-2mm}Firma del asesor: & \\
\cline{2-2} & \hspace*{1em}\asesor\hspace*{1em}
\end{tabular}




\mainmatter
\pagestyle{fime}

%%% Haz un documento para cada capítulo
The main idea behind this thesis
is to develop models and methods
for improving the service offered by car insurance agents,
helping them arrive to accident sites sooner,
and determining the number of adjusters required
to perform the service within the desired standards.

\section{Models}
%This model was tested 
%with a benchmark instance of 30 points, 
%which sets $V$ and $W$ are the same.
%For values of $p = 1,\ldots,7$
%the optimal solution was found in less than one minute.
%
%For other benchmark instances
%was impossible to find the solution
%because they have more than 300 points,
%and the proposed model grows to over one million variables.
%  For some large instances ($n = 50$ , $m = 30$)%

These test were evaluated in
a Lanix Spine BW Processor Intel Xenon,
CPU E5-2867W, 3.10 GHz.
With operative system Ubuntu 14.04.3 LTS
The models were solved with CPLEX 12.6.0.0,
and they were coded in C++.

Different random instances were generated
to validate the proposed formulations.
For $n = 100$ we created 8 instances with $m = 20,30,\ldots,100$,
and they were tested with different values of $p$ from 5 to $\frac{m}{2}$,
shown in the following table \ref{tab:tests}.
\begin{table}
  \centerign
  \label{tab:tests}
  \begin{tabular}{|c|c|c|}\hline
    n & m & p \\ \hline
    100 & 20 & 5-10 \\
    100 & 30 & 5-15 \\
    100 & 40 & 5-20 \\
    100 & 50 & 5-25 \\
    100 & 60 & 5-30 \\
    100 & 70 & 5-35 \\
    100 & 80 & 5-40 \\
    100 & 90 & 5-45 \\
    100 & 100 & 5-50 \\
    \hline
  \end{tabular}
\end{table}

Model A was solved to optimality only when $p < 25$,
in an average time of 217 seconds.
for the rest of the values of $p$, 
no factible solution was found for $65\%$ of the test cases
in less than 20 minutes,
for the rest test cases,
an average gap of $27\%$ was obtained.

Model B was evaluated
with the same combinations of $n,m,p$
and $\ell$ from $1,\ldots,p$.
For almost all cases,
model B was solved to optimality
in less than a minute for each case.
When the value of $\ell$ 
is close to $p$
model resolution takes longer
and in some cases
optimal is not achieved.
However,
it can be seen
that for different values of $\ell$
the solutions obtained are the same.


\appendix
%%% Haz un documento para cada apéndice
\chapter{Este es un apéndice}

\section{Citas bibliográficas}

En princicpio tienes total libertad de incluir tu bibliografía con el entorno {\tt thebibliography} nativo de \LaTeX{} o mediante la herramienta \textsc{Bib}\TeX. En caso de que optes por esto último (recomendado), puedes usar alguno de los archivos {\tt mighelbib.bst} o {\tt mighelnat.bst} incluidos en el paquete {\tt Tesis-FIME}, pues sus diseños están basados en el estilo bibliográfico estándar del español, además de que armoniza con el estilo de tesis provisto por {\tt fime.cls}.

El estilo bibliográfico {\tt mighelbib} es numérico, es decir cita con un número entre corchetes, por ejemplo una cita \verb+\cite{Dan82}+ genera una etiqueta del tipo [13], mientras que el estilo {\tt mighelnat} es tipo autor-año y requiere que el paquete {\tt natbib} sea cargado (sin opciones) para su correcto funcionamiento, cita con el apellido del autor y el año, por ejemplo una cita \verb+\citet{Dan82}+ genera una etiqueta del tipo Dantzig (1982), mientras que una cita \verb+\citep{Dan82}+ genera una etiqueta del tipo (Dantzig, 1982).

Como muestra del estilo, unas citas: un libro clásico de programación lineal \cite{Baz04} y un documento histórico \cite{Dan82}. Para saber un poco más del uso de \textsc{Bib}\TeX, se puede leer \cite{Mat11}.

\section{Comillas}

El objetivo de esta sección era provocar otra página para que se vea el encabezado. Pero aprovechamos para decir que la clase {\tt fime.cls} carga el paquete {\tt babel} con la opción {\tt spanish}, por lo que cambiará automáticamente los dobles signos $<<$ y $>>$ por << y >>. Estas comillas angulares son las correctas en el idioma español, y son las que se usan en la clase {\tt fime.cls}, por lo que se sugiere sean las usadas en el texto cada que quieras <<entrecomillar>> algo.



\backmatter
\pagestyle{main}

%%% Aquí va la bibliografía, puedes usar el entorno de LaTeX (thebibliography)
%%% o la herramienta BibTeX. En caso de que optes por BibTeX, puedes usar
%%% alguno de los archivos de estilo (mighelbib.bst o mighelnat.bst) incluidos
%%% en el paquete, cuyos diseños armonizan con el diseño de tesis provisto por
%%% fime.cls. Para muestra, basta un botón:
\bibliographystyle{mighelbib}
\bibliography{MiBiblio}

\label{lastpage}
%Autobiografia

\chapter*{Resumen autobiográfico}
\thispagestyle{empty}

\begin{center}
\autor

Candidato para obtener el grado de\\
\grado\\
\orientacion\bigskip

\uanl\\
\fime\bigskip

Tesis:\\
\textsc{\large\titulo}
\end{center}\bigskip

%Aquí va tu historia
Aquí va tu historia. Recuerda que debe incluir: lugar y fecha de nacimiento, nombre de los padres, escuelas y universidades en las que se graduó después de la preparatoria, títulos o grados obtenidos (no incluir los estudios que se están concluyendo), experiencia profesional y organizaciones profesionales a las que pertenece (no incluir lista de publicaciones).


\end{document}
